% GENERATED FILE. Edit canonical lessons, then run npm run generate:books.
% canonical-source-hash: fnv1a64:0ebba9bc82ad4e33
% canonical-lessons: TA-C502-anucari, TA-C502-nerunku, TA-C502-taluvu, TA-C502-varutu, TA-C502-pantam

\chapter{Drawing Close}
\label{ch:ta-drawing-close}
\hlchaptermodality{502}

\begin{chapteropening}
\textbf{By the end of this chapter:} \emph{I can talk about adjusting to others, drawing near, embracing, caressing and family bonds.}

Complete the last lesson of chapter 502: I can talk about adjusting to others, drawing near, embracing, caressing and family bonds.
\end{chapteropening}

\section[\texorpdfstring{\ta{அனுசரி} (aṉucari)}{அனுசரி (aṉucari)}]{\ta{அனுசரி} (aṉucari) — to accommodate; to adjust to}

\label{lesson:TA-C502-anucari}



\begin{quote}
Before the new one: say the Tamil for to sulk, then the Tamil for a divorce.
\end{quote}

\subsection*{You'll want to know: \ta{அனுசரி}}
\textbf{\ta{அனுசரி}} — \emph{aṉucari} — \textquotedblleft{}to accommodate; to adjust to\textquotedblright{}.

\begin{quote}
\ta{புது} \ta{வீட்டில்} \ta{எல்லோரையும்} \ta{அனுசரித்து} \ta{அவள்} \ta{வாழ்ந்தாள்}.
\end{quote}

\begin{quote}
\emph{putu vīṭṭil ellōraiyum aṉucarittu avaḷ vāḻntāḷ.}
\end{quote}

\begin{quote}
In her new home she lived in harmony with everyone.
\end{quote}

Past (I …): \ta{அனுசரித்தேன்}, \emph{aṉucarittēṉ}.

Also to observe a custom: \ta{விரதம்} \ta{அனுசரி}.

\begin{cousinweb}
\textbf{\ta{அனுசரி}} (\emph{aṉucari}) is borrowed from Sanskrit \emph{anusarati}, \textquotedblleft{}runs after, follows\textquotedblright{} (CDIAL 338), made into a Tamil verb in -\ta{இ}.
\end{cousinweb}

\begin{grammarlens}[title={What you've built}]
One more word for describing what you lived and felt.
\end{grammarlens}

\subsection*{Your turn}
\begin{itemize}
\raggedright
  \item \emph{Say it:} \emph{aṉucari}
  \item \emph{Say it:} \emph{aṉucari}, once more
  \item \emph{Say it:} \emph{putu vīṭṭil ellōraiyum aṉucarittu avaḷ vāḻntāḷ.} — In her new home she lived in harmony with everyone.
  \item \emph{Recall:} say the Tamil for to sulk, then the Tamil for a divorce, then say \emph{aṉucari} again
\end{itemize}

\begin{tcolorbox}[breakable,colback=teal!4,colframe=teal!35!black,title={Before you move on}]
What does \ta{அனுசரி} mean? (\textquotedblleft{}To accommodate; to adjust to\textquotedblright{}.) Say it once more.
\end{tcolorbox}
\section[\texorpdfstring{\ta{நெருங்கு} (neruṅku)}{நெருங்கு (neruṅku)}]{\ta{நெருங்கு} (neruṅku) — to come close; to draw near}

\label{lesson:TA-C502-nerunku}



\begin{quote}
Before the new one: say the Tamil for a divorce, then the Tamil for to accommodate.
\end{quote}

\subsection*{You'll want to know: \ta{நெருங்கு}}
\textbf{\ta{நெருங்கு}} — \emph{neruṅku} — \textquotedblleft{}to come close; to draw near\textquotedblright{}.

\begin{quote}
\ta{தேர்வு} \ta{நெருங்கியதால்} \ta{அவன்} \ta{இரவிலும்} \ta{படித்தான்}.
\end{quote}

\begin{quote}
\emph{tērvu neruṅkiyatāl avaṉ iravilum paṭittāṉ.}
\end{quote}

\begin{quote}
With the exam drawing near, he studied at night too.
\end{quote}

Past (I …): \ta{நெருங்கினேன்}, \emph{neruṅkiṉēṉ}.

\ta{நெருக்கமான} \ta{நண்பன்}, a close friend.

\begin{cousinweb}
\textbf{\ta{நெருங்கு}} (\emph{neruṅku}) is native Dravidian (DEDR 3770), in the family of \ta{நேர்} \emph{nēr}, \textquotedblleft{}to meet, to come near\textquotedblright{}; the same entry gives \ta{நெருக்கம்} \emph{nerukkam}, \textquotedblleft{}closeness\textquotedblright{}.
\end{cousinweb}

\begin{grammarlens}[title={What you've built}]
One more word for describing what you lived and felt.
\end{grammarlens}

\subsection*{Your turn}
\begin{itemize}
\raggedright
  \item \emph{Say it:} \emph{neruṅku}
  \item \emph{Say it:} \emph{neruṅku}, once more
  \item \emph{Say it:} \emph{tērvu neruṅkiyatāl avaṉ iravilum paṭittāṉ.} — With the exam drawing near, he studied at night too.
  \item \emph{Recall:} say the Tamil for a divorce, then the Tamil for to accommodate, then say \emph{neruṅku} again
\end{itemize}

\begin{tcolorbox}[breakable,colback=teal!4,colframe=teal!35!black,title={Before you move on}]
What does \ta{நெருங்கு} mean? (\textquotedblleft{}To come close; to draw near\textquotedblright{}.) Say it once more.
\end{tcolorbox}
\section[\texorpdfstring{\ta{தழுவு} (taḻuvu)}{தழுவு (taḻuvu)}]{\ta{தழுவு} (taḻuvu) — to embrace}

\label{lesson:TA-C502-taluvu}



\begin{quote}
Before the new one: say the Tamil for to accommodate, then the Tamil for to come close.
\end{quote}

\subsection*{You'll want to know: \ta{தழுவு}}
\textbf{\ta{தழுவு}} — \emph{taḻuvu} — \textquotedblleft{}to embrace\textquotedblright{}.

\begin{quote}
\ta{பல} \ta{வருடம்} \ta{கழித்து} \ta{வந்த} \ta{மகனை} \ta{அம்மா} \ta{தழுவினாள்}.
\end{quote}

\begin{quote}
\emph{pala varuṭam kaḻittu vanta makaṉai ammā taḻuviṉāḷ.}
\end{quote}

\begin{quote}
The mother embraced the son who came home after many years.
\end{quote}

Past (I …): \ta{தழுவினேன்}, \emph{taḻuviṉēṉ}.

Also to adopt an idea or a style: \ta{மேற்கத்திய} \ta{முறையைத்} \ta{தழுவு}.

\begin{cousinweb}
\textbf{\ta{தழுவு}} (\emph{taḻuvu}) is native Dravidian (DEDR 3116), to clasp and to hug.
\end{cousinweb}

\begin{grammarlens}[title={What you've built}]
One more word for describing what you lived and felt.
\end{grammarlens}

\subsection*{Your turn}
\begin{itemize}
\raggedright
  \item \emph{Say it:} \emph{taḻuvu}
  \item \emph{Say it:} \emph{taḻuvu}, once more
  \item \emph{Say it:} \emph{pala varuṭam kaḻittu vanta makaṉai ammā taḻuviṉāḷ.} — The mother embraced the son who came home after many years.
  \item \emph{Recall:} say the Tamil for to accommodate, then the Tamil for to come close, then say \emph{taḻuvu} again
\end{itemize}

\begin{tcolorbox}[breakable,colback=teal!4,colframe=teal!35!black,title={Before you move on}]
What does \ta{தழுவு} mean? (\textquotedblleft{}To embrace\textquotedblright{}.) Say it once more.
\end{tcolorbox}
\section[\texorpdfstring{\ta{வருடு} (varuṭu)}{வருடு (varuṭu)}]{\ta{வருடு} (varuṭu) — to stroke; to caress}

\label{lesson:TA-C502-varutu}



\begin{quote}
Before the new one: say the Tamil for to come close, then the Tamil for to embrace.
\end{quote}

\subsection*{You'll want to know: \ta{வருடு}}
\textbf{\ta{வருடு}} — \emph{varuṭu} — \textquotedblleft{}to stroke; to caress\textquotedblright{}.

\begin{quote}
\ta{அம்மா} \ta{என்} \ta{தலையை} \ta{அன்போடு} \ta{வருடினாள்}.
\end{quote}

\begin{quote}
\emph{ammā eṉ talaiyai aṉpōṭu varuṭiṉāḷ.}
\end{quote}

\begin{quote}
Mother stroked my head lovingly.
\end{quote}

Past (I …): \ta{வருடினேன்}, \emph{varuṭiṉēṉ}.

Gently, with the hand.

\begin{cousinweb}
\textbf{\ta{வருடு}} (\emph{varuṭu}) is a native Tamil verb whose origin is not settled; DEDR has no entry for it that we can point to.
\end{cousinweb}

\begin{grammarlens}[title={What you've built}]
One more word for describing what you lived and felt.
\end{grammarlens}

\subsection*{Your turn}
\begin{itemize}
\raggedright
  \item \emph{Say it:} \emph{varuṭu}
  \item \emph{Say it:} \emph{varuṭu}, once more
  \item \emph{Say it:} \emph{ammā eṉ talaiyai aṉpōṭu varuṭiṉāḷ.} — Mother stroked my head lovingly.
  \item \emph{Recall:} say the Tamil for to come close, then the Tamil for to embrace, then say \emph{varuṭu} again
\end{itemize}

\begin{tcolorbox}[breakable,colback=teal!4,colframe=teal!35!black,title={Before you move on}]
What does \ta{வருடு} mean? (\textquotedblleft{}To stroke; to caress\textquotedblright{}.) Say it once more.
\end{tcolorbox}
\section[\texorpdfstring{\ta{பந்தம்} (pantam)}{பந்தம் (pantam)}]{\ta{பந்தம்} (pantam) — a bond; a tie}

\label{lesson:TA-C502-pantam}



\begin{quote}
Before the new one: say the Tamil for to embrace, then the Tamil for to stroke.
\end{quote}

\subsection*{You'll want to know: \ta{பந்தம்}}
\textbf{\ta{பந்தம்}} — \emph{pantam} — \textquotedblleft{}a bond; a tie\textquotedblright{}.

\begin{quote}
\ta{திருமணத்துக்கு} \ta{எல்லாச்} \ta{சொந்த} \ta{பந்தங்களும்} \ta{வந்தார்கள்}.
\end{quote}

\begin{quote}
\emph{tirumaṇattukku ellāc conta pantaṅkaḷum vantārkaḷ.}
\end{quote}

\begin{quote}
All our relatives and kin came for the wedding.
\end{quote}

\ta{சொந்த} \ta{பந்தம்}, kith and kin.

\begin{cousinweb}
\textbf{\ta{பந்தம்}} (\emph{pantam}) is borrowed from Sanskrit \emph{bandha}, \textquotedblleft{}a bond\textquotedblright{} (CDIAL 9136). DEDR lists a native \ta{பந்தம்} \textquotedblleft{}torch\textquotedblright{} (3920), which is a different word.
\end{cousinweb}

\begin{grammarlens}[title={What you've built}]
One more word for describing what you lived and felt.
\end{grammarlens}

\subsection*{Your turn}
\begin{itemize}
\raggedright
  \item \emph{Say it:} \emph{pantam}
  \item \emph{Say it:} \emph{pantam}, once more
  \item \emph{Say it:} \emph{tirumaṇattukku ellāc conta pantaṅkaḷum vantārkaḷ.} — All our relatives and kin came for the wedding.
  \item \emph{Recall:} say the Tamil for to embrace, then the Tamil for to stroke, then say \emph{pantam} again
\end{itemize}

\begin{tcolorbox}[breakable,colback=teal!4,colframe=teal!35!black,title={Before you move on}]
What does \ta{பந்தம்} mean? (\textquotedblleft{}A bond; a tie\textquotedblright{}.) Say it once more.
\end{tcolorbox}
